% GENERATED FILE. Edit canonical lessons, then run npm run generate:books.
% canonical-source-hash: fnv1a64:a81d64c231db1a95
% canonical-lessons: HI-C332-parampara, HI-C332-rivaj, HI-C332-utsav, HI-C332-alav, HI-C332-putla

\chapter{Old Customs}
\label{ch:hi-old-customs}
\hlchaptermodality{332}

\begin{chapteropening}
\textbf{By the end of this chapter:} \emph{I can explain an old custom and how it is kept.}

Complete the last lesson of chapter 332: I can explain an old custom and how it is kept.
\end{chapteropening}

\section[\texorpdfstring{\hi{परंपरा} (paramparā)}{परंपरा (paramparā)}]{\hi{परंपरा} (paramparā) — a tradition}

\label{lesson:HI-C332-parampara}



\begin{quote}
Before the new one: say the Hindi for a blessed offering (of food), then the Hindi for a festive dish.
\end{quote}

\subsection*{You'll want to know: \hi{परंपरा}}
\textbf{\hi{परंपरा}} — \emph{paramparā} — \textquotedblleft{}a tradition\textquotedblright{}.

\begin{quote}
\hi{यह} \hi{परंपरा} \hi{सौ} \hi{साल} \hi{पुरानी} \hi{है।}
\end{quote}

\begin{quote}
\emph{yah paramparā sau sāl purānī hai.}
\end{quote}

\begin{quote}
This tradition is a hundred years old.
\end{quote}

Feminine.

\begin{cousinweb}
\textbf{\hi{परंपरा}} (\emph{paramparā}) is Sanskrit \emph{paramparā}, \textquotedblleft{}an unbroken series, one after another\textquotedblright{} (Monier-Williams): what passes from each generation to the next.
\end{cousinweb}

\begin{grammarlens}[title={What you've built}]
One more word for telling what happened.
\end{grammarlens}

\subsection*{Your turn}
\begin{itemize}
\raggedright
  \item \emph{Say it:} \emph{paramparā}
  \item \emph{Say it:} \emph{paramparā}, once more
  \item \emph{Say it:} \emph{yah paramparā sau sāl purānī hai.} — This tradition is a hundred years old.
  \item \emph{Recall:} say the Hindi for a blessed offering (of food), then the Hindi for a festive dish, then say \emph{paramparā} again
\end{itemize}

\begin{tcolorbox}[breakable,colback=teal!4,colframe=teal!35!black,title={Before you move on}]
What does \hi{परंपरा} mean? (\textquotedblleft{}A tradition\textquotedblright{}.) Say it once more.
\end{tcolorbox}
\section[\texorpdfstring{\hi{रिवाज} (rivāj)}{रिवाज (rivāj)}]{\hi{रिवाज} (rivāj) — a custom}

\label{lesson:HI-C332-rivaj}



\begin{quote}
Before the new one: say the Hindi for a festive dish, then the Hindi for a tradition.
\end{quote}

\subsection*{You'll want to know: \hi{रिवाज}}
\textbf{\hi{रिवाज}} — \emph{rivāj} — \textquotedblleft{}a custom\textquotedblright{}.

\begin{quote}
\hi{हमारे} \hi{गाँव} \hi{में} \hi{त्योहार} \hi{पर} \hi{नए} \hi{कपड़े} \hi{पहनने} \hi{का} \hi{रिवाज} \hi{है।}
\end{quote}

\begin{quote}
\emph{hamāre gā̃v meṁ tyohār par nae kapṛe pahanne kā rivāj hai.}
\end{quote}

\begin{quote}
In our village it is the custom to wear new clothes for a festival.
\end{quote}

Masculine.

\begin{cousinweb}
\textbf{\hi{रिवाज}} (\emph{rivāj}) is usually traced, through Persian, to Arabic \emph{rawāj}, \textquotedblleft{}currency, being in use\textquotedblright{}: what is current in a place. \textbf{\hi{का} \hi{रिवाज} \hi{है}} is it is the custom to.
\end{cousinweb}

\begin{grammarlens}[title={What you've built}]
One more word for telling what happened.
\end{grammarlens}

\subsection*{Your turn}
\begin{itemize}
\raggedright
  \item \emph{Say it:} \emph{rivāj}
  \item \emph{Say it:} \emph{rivāj}, once more
  \item \emph{Say it:} \emph{hamāre gā̃v meṁ tyohār par nae kapṛe pahanne kā rivāj hai.} — In our village it is the custom to wear new clothes for a festival.
  \item \emph{Recall:} say the Hindi for a festive dish, then the Hindi for a tradition, then say \emph{rivāj} again
\end{itemize}

\begin{tcolorbox}[breakable,colback=teal!4,colframe=teal!35!black,title={Before you move on}]
What does \hi{रिवाज} mean? (\textquotedblleft{}A custom\textquotedblright{}.) Say it once more.
\end{tcolorbox}
\section[\texorpdfstring{\hi{उत्सव} (utsav)}{उत्सव (utsav)}]{\hi{उत्सव} (utsav) — a celebration; a festival}

\label{lesson:HI-C332-utsav}



\begin{quote}
Before the new one: say the Hindi for a tradition, then the Hindi for a custom.
\end{quote}

\subsection*{You'll want to know: \hi{उत्सव}}
\textbf{\hi{उत्सव}} — \emph{utsav} — \textquotedblleft{}a celebration; a festival\textquotedblright{}.

\begin{quote}
\hi{गाँव} \hi{का} \hi{उत्सव} \hi{तीन} \hi{दिन} \hi{चला।}
\end{quote}

\begin{quote}
\emph{gā̃v kā utsav tīn din calā.}
\end{quote}

\begin{quote}
The village celebration went on for three days.
\end{quote}

Masculine.

\begin{cousinweb}
\textbf{\hi{उत्सव}} (\emph{utsav}) is Sanskrit \emph{utsava}, \textquotedblleft{}a festival, a joyful occasion\textquotedblright{} (Monier-Williams), the formal word beside \textbf{\hi{त्योहार}}.
\end{cousinweb}

\begin{grammarlens}[title={What you've built}]
One more word for telling what happened.
\end{grammarlens}

\subsection*{Your turn}
\begin{itemize}
\raggedright
  \item \emph{Say it:} \emph{utsav}
  \item \emph{Say it:} \emph{utsav}, once more
  \item \emph{Say it:} \emph{gā̃v kā utsav tīn din calā.} — The village celebration went on for three days.
  \item \emph{Recall:} say the Hindi for a tradition, then the Hindi for a custom, then say \emph{utsav} again
\end{itemize}

\begin{tcolorbox}[breakable,colback=teal!4,colframe=teal!35!black,title={Before you move on}]
What does \hi{उत्सव} mean? (\textquotedblleft{}A celebration; a festival\textquotedblright{}.) Say it once more.
\end{tcolorbox}
\section[\texorpdfstring{\hi{अलाव} (alāv)}{अलाव (alāv)}]{\hi{अलाव} (alāv) — a bonfire}

\label{lesson:HI-C332-alav}



\begin{quote}
Before the new one: say the Hindi for a custom, then the Hindi for a celebration.
\end{quote}

\subsection*{You'll want to know: \hi{अलाव}}
\textbf{\hi{अलाव}} — \emph{alāv} — \textquotedblleft{}a bonfire\textquotedblright{}.

\begin{quote}
\hi{ठंड} \hi{की} \hi{रात} \hi{में} \hi{लोग} \hi{अलाव} \hi{के} \hi{पास} \hi{बैठे।}
\end{quote}

\begin{quote}
\emph{ṭhaṇḍ kī rāt meṁ log alāv ke pās baiṭhe.}
\end{quote}

\begin{quote}
On the cold night people sat around the bonfire.
\end{quote}

Masculine.

\begin{cousinweb}
\textbf{\hi{अलाव}} (\emph{alāv}) goes back to Sanskrit \emph{alāta}, \textquotedblleft{}a firebrand, a burning stick\textquotedblright{} (CDIAL 710): the open fire people sit around in winter.
\end{cousinweb}

\begin{grammarlens}[title={What you've built}]
One more word for telling what happened.
\end{grammarlens}

\subsection*{Your turn}
\begin{itemize}
\raggedright
  \item \emph{Say it:} \emph{alāv}
  \item \emph{Say it:} \emph{alāv}, once more
  \item \emph{Say it:} \emph{ṭhaṇḍ kī rāt meṁ log alāv ke pās baiṭhe.} — On the cold night people sat around the bonfire.
  \item \emph{Recall:} say the Hindi for a custom, then the Hindi for a celebration, then say \emph{alāv} again
\end{itemize}

\begin{tcolorbox}[breakable,colback=teal!4,colframe=teal!35!black,title={Before you move on}]
What does \hi{अलाव} mean? (\textquotedblleft{}A bonfire\textquotedblright{}.) Say it once more.
\end{tcolorbox}
\section[\texorpdfstring{\hi{पुतला} (putlā)}{पुतला (putlā)}]{\hi{पुतला} (putlā) — an effigy; a large doll}

\label{lesson:HI-C332-putla}



\begin{quote}
Before the new one: say the Hindi for a celebration, then the Hindi for a bonfire.
\end{quote}

\subsection*{You'll want to know: \hi{पुतला}}
\textbf{\hi{पुतला}} — \emph{putlā} — \textquotedblleft{}an effigy; a large doll\textquotedblright{}.

\begin{quote}
\hi{त्योहार} \hi{की} \hi{रात} \hi{लोगों} \hi{ने} \hi{एक} \hi{बड़ा} \hi{पुतला} \hi{जलाया।}
\end{quote}

\begin{quote}
\emph{tyohār kī rāt logoṁ ne ek baṛā putlā jalāyā.}
\end{quote}

\begin{quote}
On the festival night people burned a huge effigy.
\end{quote}

Masculine.

\begin{cousinweb}
is inherited from \emph{*putrala}, a form CDIAL reconstructs (CDIAL 8269), with the sense of a doll. A huge \textbf{\hi{पुतला}} of a demon is burned at some festivals.
\end{cousinweb}

\begin{grammarlens}[title={What you've built}]
One more word for telling what happened.
\end{grammarlens}

\subsection*{Your turn}
\begin{itemize}
\raggedright
  \item \emph{Say it:} \emph{putlā}
  \item \emph{Say it:} \emph{putlā}, once more
  \item \emph{Say it:} \emph{tyohār kī rāt logoṁ ne ek baṛā putlā jalāyā.} — On the festival night people burned a huge effigy.
  \item \emph{Recall:} say the Hindi for a celebration, then the Hindi for a bonfire, then say \emph{putlā} again
\end{itemize}

\begin{tcolorbox}[breakable,colback=teal!4,colframe=teal!35!black,title={Before you move on}]
What does \hi{पुतला} mean? (\textquotedblleft{}An effigy; a large doll\textquotedblright{}.) Say it once more.
\end{tcolorbox}
